\section{Discussion}
\label{sec:Diskussion}
An important constriction for the conduction of the experiment was to obtain a total particle count of $N \geq 1112$ to keep the relative error low.
The data in the appendix shows that this could not be realized within a reasonable time frame for neither of the measurements, especially for large scattering angles.
It therefore can be assumed that the error in the count rate causes relevant deviations from the theoretically expected scattering distribution.
%
\subsection{Discussion of the detector signal pulses}
In the two signal pulses in \autoref{fig:pulses} the positive influence of the amplification is clearly visible.
The signal-to-noise ratio gets smaller and the triggering with the oscilloscope becomes easyer and less faulty.
Nevertheless, a direct comparison of the two pulses is not possible because of the differently chosen timestep settings.
While \autoref{subfig:pre} shows the noise signal in detail and underlines the difficulty of triggering to the real signal pulse, the characteristic falling flank of the signal cannot be observed.
In \autoref{subfig:after}, this flank can be seen clearly.
Although, the timestep setting does not show the characteristics of the noise.
\newline\noindent
For better comparison a setting should be chosen in which both signal pulses fill out the oscilloscope screen as best as possible.
Since the pulse length is not amplified, this should be possible by keeping the same timestep setting and only adjusting the voltage setting to account for the amplification factor.
%
\subsection{Usage of the calibration measurement}
The measurement of the $\alpha$-particle beam profile shows that the alignment of the beam and the detector is off by $\theta = \qty{1(0.1)}{\degree}$.
This could be due to the center-point being wrongly defined before or by the radiation source being placed off-axis.
In both cases the whole beam profile would be a mirror-symmetric gaussian curve with the maximum at the shifted angle.
\newline\noindent
In \autoref{fig:calibration} it can be seen that the beam profile is not symmetric around the red dashed line, which marks the shifted maximum.
An explanation for that behaviour could be the statistical error of the count rate because of the randomness of the atomic decay.
Also, a slight misplacement between the measurement points could cause the deviation.
Before starting the measurement, the radioactive source had to be taped down inside the vacuum chamber.
Otherwise, the vibrations of the vacuum pump cause the source to move uncontrolled.
It is therefore possible, though it cannot be quantified, that the tension of the duct tape causes slight movements. 
%
\subsection{Angular dependence of scattering for all used materials}
The Rutherford scattering formula \autoref{eqn:rutherford} states that the theoretical value for the exponent on the sine function is $b_{\mathrm{theo}} = \qty{4}{}$.
This value can be compared to the fit parameters for the measurment with gold, bismuth, and aluminum foil by using the relative deviation formula from \autoref{eqn:relabw}.
The parameter values and their relative deviations are listed in \autoref{tab:relabw}.
\begin{table}
  \centering
  \caption{Relative deviations between experimental and theoretical value for the exponent $b$ of the sine in the Rutherford scattering formula~\ref{eqn:rutherford}.}
  \label{tab:relabw}
  \sisetup{table-format=1.1, per-mode=reciprocal}
  \begin{tblr}{
      colspec = {S[table-format=3.0] },
      row{1} = {guard, mode=math},
      vline{3,5} = {2}{-}{text=\clap{$\pm$}},
    }
    \toprule
    \mathrm{Material} & \SetCell[c=2]{c} b_{\mathrm{exp}} & & \SetCell[c=2]{c} \increment b_{\mathrm{rel}} \mathbin{/} \unit{\percent} & \\
    \midrule
        \text{Gold}     & 4.05 & 0.09 & 1.25 & 2.25 \\
        \text{Bismuth}  & 4.29 & 0.04 & 7.25 & 1.0  \\
        \text{Aluminum} & 4.31 & 0.06 & 7.75 & 1.50 \\
    \bottomrule
  \end{tblr}
\end{table}
\FloatBarrier
\noindent
The small deviations of under $\qty{10}{\percent}$ for all three measurements show that the experiment reproduced the theoretical expectations by reasonable accuracy.
Still, the error of the experimental values is large in comparison to the relative error of the nominal value.
This is best visible for the measurement with gold foil, in which the error on the relative deviation is nearly double the nominal deviation.
For the other two experiments this effect is smaller, because the relative deviation is larger in the first place.
Therefore, the valid quantity is the nominal value of the relative deviation which is small in all three cases.
The large errors can be explained by the already discussed large statistical errors of the total particle counts due to time restrictions and low count rates.
\newline\noindent
Not all data points have been included into the fit functions of the three measurements.
In case of the gold and bismuth measurement this is due to the small scattering angle values causing the exponent $b$ to vary largly from the expected value.
There also is a visible difference in the proportionality for the large and small scattering angle values.
This is most significant in \autoref{fig:Au2}, where the values up to $\theta = \qty{6(0.1)}{\degree}$ follow a less steep curve than the data for the larger angles.
For the bismuth measurement this effect also occurs, but less strong.
\newline\noindent
The data for the aluminum measurement differs from that behaviour.
\autoref{fig:Al3} shows, that only the third and fourth value vary from the other data points.
Opposite to the other two measurements, the small and large angle data correspond to the same fit function.
Including the third and fourth value led to the fit not converging, even if the smallest angle values were left out.
Therefore, the third and fourth value are probably faulty data points.
\newline\noindent
Furthermore, the aluminum foil is the thickest of the three used foils.
Therefore, multiple scattering could lead to measurement errors such as those which were observed.
%
\subsection{Discussion of the Z-dependence}
For checking the $Z$-dependence a large scattering angle is chosen, because for large angles the count rate only differs slightly for small variations of the scattering angle.
Choosing a larger scattering angle than $\theta = \qty{9(0.1)}{\degree}$ would therefore be beneficial.
The aluminum measurement is the constraint in this case.
\newline\noindent
The plotted data in \autoref{fig:Z_dependence} shows no clear proportionality to $Z^2$.
Although it is possible to fit the data points with a quadratic equation, the fitted graph does not intersect with the last two data points or their errorbars.
The fit parameters also show, that the fit function does not fit the data.
For the parameter $B$ the ratio between nominal value and standard deviation is $\frac{B}{\sigma B} = \qty{97.5}{\percent}$.
\newline\noindent
Therefore, the theoretically proposed $Z$-dependence cannot be confirmed or denied, because the data is not sufficient and probably faulty.
An improved experiment could include more materials to provide more data along the expected curve.
Also, the measurement should be focused on one large scattering angle or multiple angles, for which the measurement time is increased to one or more hours.